\section{Design of Cstamp-PWAL}
\label{sec:design}

This section describes the design of Cstamp-PWAL, focusing on three key ideas:
reusing the commit timestamp as the logical recovery order, computing dependency
frontiers for safe acknowledgment, and separating the worker, flusher, and
committer into a pipeline.

\subsection{Reusing Commit Timestamp as Logical LSN}

In timestamp-based engines like ERMIA, the commit timestamp already defines the
serialization order of all committed transactions.  Therefore, we can eliminate a
separate global WAL LSN and instead use the commit timestamp as the logical
recovery order.

\begin{definition}[Logical LSN]
The logical LSN of a transaction is its commit timestamp assigned by the
concurrency control layer at validation time.
\end{definition}

This simplification has two benefits:

\begin{enumerate}
  \item \textbf{Eliminates a hot-path allocation}: No need to allocate a separate
        monotonic counter for WAL LSNs.
  \item \textbf{Tightens CC/logging coupling}: The same timestamp determines both
        serializability and recovery order, reducing the state space and potential
        for errors.
\end{enumerate}

\subsection{Dependency Frontier and Safe Acknowledgment}

Instead of waiting for a global durable prefix, Cstamp-PWAL uses the read-write
dependency frontier to determine when a transaction can be acknowledged.

\begin{definition}[Dependency Frontier]
The dependency frontier of a transaction $T$ is the set of all transactions that
$T$ transitively depends on via read-write relationships, plus $T$ itself.
\end{definition}

The system maintains this frontier by:

\begin{enumerate}

\item \textbf{Detecting dependencies at validation}: During transaction
validation, the system records which other committed transactions wrote to data
items that this transaction read.

\item \textbf{Computing transitive closure}: The frontier includes direct
dependencies and all transitive ancestors.

\item \textbf{Acknowledging only when frontier is durable}: A transaction is
acknowledged as durable only when all transactions in its frontier have durable
log records.

\end{enumerate}

This condition is safe because:
- If $T$ depends on $U$, then $U$ is in $T$'s frontier.
- If $U$ is in the frontier and becomes durable, the effects $T$ depended on are
  recoverable.
- Therefore, acknowledging $T$ after its frontier is durable ensures recovery
  consistency.

On sparse-dependency workloads, the frontier is small, so Cstamp-PWAL can
acknowledge transactions quickly without waiting for unrelated shards.  On dense-
dependency workloads, the frontier grows and asymptotically approaches a global
durable prefix, which is correct though conservative.

\subsection{Worker / Flusher / Committer Pipeline}

Cstamp-PWAL separates durability concerns across three types of threads to
avoid blocking worker threads on slow operations:

\begin{enumerate}

\item \textbf{Worker threads}: Execute transactions, validate them, assign commit
timestamps, and enqueue log records to shards.  They never block on
\texttt{fdatasync} or on durable-prefix waits.

\item \textbf{Flusher threads}: Read log records from shard queues, batch them,
and persist them to stable storage via \texttt{fdatasync}.  Each shard has its
own flusher, allowing parallel flushing.

\item \textbf{Committer thread}: Observes updates to each shard's durable frontier
and determines which transactions' dependency conditions are satisfied.
Transactions whose entire frontier is durable are acknowledged and removed from
pending.

\end{enumerate}

This pipelining allows:
- Workers to continue executing while flushers are persisting, without blocking.
- Committers to work on a potentially stale view of durable frontiers, as long as
  they eventually become correct.
- The system to batch flushes for efficiency without workers paying the latency
  cost directly.

\subsection{Integration with ERMIA}

Cstamp-PWAL integrates with ERMIA as follows:

\begin{itemize}

\item \textbf{Dependency detection}: We leverage ERMIA's validation logic.  When a
transaction validates, it already records which earlier committed versions it read
from.  We extract this information to build the read-write dependency relation.

\item \textbf{Commit timestamp}: ERMIA's commit timestamp, already assigned at
validation, becomes the logical recovery order.  No separate allocation is needed.

\item \textbf{Version management}: ERMIA's multi-version storage and timestamp-
based version selection remain unchanged.  Cstamp-PWAL operates only on the
durability layer.

\end{itemize}

The coupling is minimal and non-invasive: Cstamp-PWAL sits as a layer between
transaction validation and the traditional WAL system, observing dependencies
and managing durable acknowledgment.

